\begin{summarybox}
\textbf{C013｜长轴端点切线相交：一眼识别、三式通杀}

\textbf{识别图形}：椭圆 $x^2/a^2+y^2/b^2=1$；
两条固定端点切线 $x=\pm a$；第三条为椭圆动切线。
注意切点 $P\ne(\pm a,0)$。

\textbf{记住三式}（$d_1,d_2$ 为两条竖段长）：
\[
\boxed{d_1d_2=b^2},\qquad
\boxed{S=a(d_1+d_2)\ge2ab},\qquad
\boxed{S_{\min}=2ab}.
\]
等号当且仅当 $d_1=d_2=b$，
切点为短轴端点 $P=(0,\pm b)$；面积无最大值。

\textbf{最快证明（斜截式）}：切线 $y=kx+m$ 满足
\[
m^2-a^2k^2=b^2;\qquad
 y_L=m-ak,\quad y_R=m+ak.
\]
于是 $y_Ly_R=b^2>0$，交点同侧，
\[
S=a|y_L+y_R|=2a|m|
=2a\sqrt{a^2k^2+b^2}\ge2ab.
\]

\textbf{最直观证明（圆）}：先取半径 $b$ 的圆。
同一点引两切线，切线段相等；对横跨圆的整条切线用勾股定理，
推出两端竖段高度积为 $b^2$；把圆沿水平方向拉长 $a/b$ 倍，
得到椭圆，竖段长度不变，故定值仍为 $b^2$。

\textbf{第二条快通道（已知切点）}：
\[
\boxed{S=\frac{2ab^2}{|y_0|}},\qquad 0<|y_0|\le b.
\]
短轴端点面积最小；靠近长轴端点面积无限增大。

\textbf{关联知识}：圆的切线长定理、勾股定理、仿射伸缩、
平方差公式、基本不等式、辅助角公式、支持线与切线判别式。

\textbf{防错一句话}：这是 $x=\pm a$ 两条\textbf{端点切线}的特殊定值，
不是随便两条竖线、更不是两条\textbf{准线}的普遍定值。
\end{summarybox}
